% =============================================================================
% Section 02 — Unity of Identity
% Book I — Ontology of Viability
% Chapter 03 — Unity and Plurality
%
% Purpose:
%   To examine the ontological meaning and scope of the unity
%   of identity without reducing identity to singularity,
%   numerical indivisibility, unity of source, composition,
%   or invariance.
%
% Conceptual Policy:
%   1. Unity of identity must be distinguished from the unity
%      of an aggregate and from the unity of a source.
%   2. Unity must not be equated with singularity.
%   3. The unity of identity must not be derived from the
%      invariance of a configuration or set of determinations.
%   4. The possibility of plurality of identities must remain
%      open and must not be excluded by definition.
%   5. No universal assumption that every existent possesses
%      exactly one identity is adopted.
%
% =============================================================================

\section{Unity of Identity}
\label{sec:unity-of-identity}

% -----------------------------------------------------------------------------
\subsection{The Meaning of Unity in the Context of Identity}

The question of unity acquires a distinct ontological significance
when it is considered in relation to identity.

An existent may be described as unified in different respects.
Its determinations may form an organized configuration;
its components may constitute an aggregate;
its manifestations may be related to a common source;
or its existence may be considered under a single identity.
These descriptions concern different possible forms of unity
and cannot be treated as interchangeable.

The present section is concerned specifically with the
\emph{unity of identity}.
It therefore asks whether, and in what sense, identity may
constitute a unity, rather than merely whether an existent
possesses an organized structure or an identifiable source.

Unity of identity is not simply the observation that several
determinations belong to one existent.
Nor is it equivalent to the fact that several components
participate in one composition.
The ontological question concerns the unity attributable
to identity itself.

This distinction is necessary because the unity of an existent
in one respect does not automatically establish its unity
in every other respect.

A composition may be unified while the question of its
identity remains unresolved.
A plurality of manifestations may be related without
thereby determining whether they belong to one identity
or to several.
Likewise, a common source may provide a basis for examining
unity without proving that every manifestation associated
with that source possesses the same identity.

Unity of identity must therefore be examined as a distinct
ontological determination.

% -----------------------------------------------------------------------------
\subsection{Unity Is Not an Additional Property of Identity}

Unity should not be understood as an independent property
that is externally attached to an already constituted identity.

If identity is ontologically inseparable from being,
then unity cannot be treated as a merely supplementary
attribute that is added to identity after its constitution.

Such an interpretation would presuppose that identity
is first constituted independently and subsequently
receives unity as an additional determination.
This would introduce an artificial separation between
identity and a possible mode of its ontological unity.

However, recognizing this inseparability does not establish
that every identity must be singular in every possible respect.
Nor does it demonstrate that the unity of identity can be
adequately described through a single universal formula.

The claim that unity is not externally added to identity
must be distinguished from the stronger claim that every
existent necessarily possesses exactly one identity.

The former concerns the ontological relation between
identity and unity.
The latter concerns the number of identities that may
belong to an existent.

These are different questions and require different arguments.

Accordingly, unity of identity is considered here as an
ontological problem concerning the mode in which identity
is unified, rather than as a supplementary attribute
or a predetermined numerical condition.

% -----------------------------------------------------------------------------
\subsection{Unity of Identity and Numerical Singularity}

Unity of identity must be distinguished from numerical
singularity.

Numerical singularity concerns whether something is one
in number rather than more than one.
Unity, however, may concern the ontological coherence
or unity of identity without necessarily settling
the numerical question.

An existent may be numerically one under a particular
description while the ontological status of its identity
remains open.

Conversely, the possibility that an existent involves
more than one identity cannot be excluded merely because
it is treated as one object, one organism, one system,
or one material aggregate in a particular domain.

The numerical designation of an existent does not,
by itself, determine the number or mode of its identities.

This distinction prevents a transition from numerical
description to ontological conclusion without sufficient
argument.

To call something one does not automatically establish
that its identity is singular.
Likewise, to identify several components or manifestations
does not automatically establish that several identities
are present.

The relation between unity and singularity is examined
more directly in Section~\ref{sec:non-equivalence-of-unity-and-singularity}.
For the present inquiry, it is sufficient to establish
that numerical singularity cannot serve as an unquestioned
definition or proof of unity of identity.

% -----------------------------------------------------------------------------
\subsection{Unity of Identity and Unity of Source}

The unity of identity must also be distinguished from
the unity of source.

Several manifestations may be associated with one source.
This relation may provide an important basis for examining
their ontological connection.
Nevertheless, the unity of their source does not
automatically determine the unity of their identities.

The inference from one source to one identity requires
additional ontological justification.

A common source may be relevant to the constitution,
generation, or manifestation of several entities.
Yet the relation between source and identity cannot
be settled merely by counting sources or manifestations.

It is therefore necessary to distinguish at least
two questions:

\begin{enumerate}
    \item Whether several manifestations are related
    to one source.

    \item Whether those manifestations belong to one
    identity or involve a plurality of identities.
\end{enumerate}

The first question concerns the ontological relation
of manifestations to their source.
The second concerns the unity or plurality of identity.

Neither question can be substituted for the other.

This distinction also prevents the reduction of identity
to origin.
An existent cannot be assigned an identity solely by
establishing where it originates, unless the ontological
relevance of that origin to identity has independently
been demonstrated.

The distinction between unity of source and unity
of aggregate is developed in
Section~\ref{sec:unity-of-source-and-aggregate}.

% -----------------------------------------------------------------------------
\subsection{Unity of Identity and Composition}

The unity of identity must not be reduced to the unity
of composition.

Composition concerns the relation among components
that constitute an organized whole.
It may involve material, structural, functional,
or other forms of organization.

Identity, however, is not exhausted by the arrangement
or aggregation of components.

A composed existent may exhibit a high degree of
structural integration without this fact alone establishing
the ontological unity of its identity.

Likewise, changes in composition do not necessarily
establish the emergence of a new identity.
Nor does the preservation of a composition necessarily
prove the persistence of one identity.

The relation between identity and composition therefore
requires an analysis that does not assume either that
composition determines identity or that identity
is independent of every compositional condition.

The relevant question is not whether composition
is ontologically irrelevant.
Rather, it is whether and under what conditions
compositional relations contribute to the unity
of identity.

This question is treated separately in
Section~\ref{sec:plurality-and-composition}.

% -----------------------------------------------------------------------------
\subsection{Unity of Identity and Invariance}

Unity of identity must also be distinguished from
the preservation of invariant determinations.

An existent may preserve certain determinations
under transformation.
Such preservation may establish an invariant relation
within a specified domain and with respect to
a specified transformation.

However, the presence of invariants does not by itself
establish that the existent possesses one unified identity.

An invariant may be shared by distinct existents.
Different identities may preserve similar or identical
determinations.
Conversely, an existent may undergo transformations
in which previously identified invariants cease to hold
without this fact alone demonstrating the disappearance
of its identity.

The unity of identity therefore cannot be inferred
solely from the persistence of invariant features.

This does not mean that invariance is irrelevant
to the examination of identity.
It means that invariance and unity of identity
answer different ontological questions.

Invariance concerns what remains preserved under
a specified transformation.
Unity of identity concerns the ontological unity
attributable to identity.

Their possible relations must be established
without reducing one to the other.

The distinction developed in Chapter~\ref{ch:invariance}
is therefore presupposed here as a conceptual constraint,
not as a complete solution to the problem of unity.

% -----------------------------------------------------------------------------
\subsection{Unity and the Plurality of Determinations}

An existent may involve a plurality of determinations.
These determinations may differ in kind, mode,
relation, or ontological role.

The presence of such plurality does not automatically
negate the possibility of unity of identity.

A plurality of determinations may be compatible
with the unity of an identity.
However, this compatibility must not be interpreted
as proof that all plural determinations necessarily
belong to one identity.

The question concerns the ontological relation
between determinations and identity.

It is therefore necessary to avoid two opposite reductions.

The first reduction assumes that plurality of
determinations necessarily implies plurality
of identities.

The second assumes that the unity of an existent
necessarily absorbs every plurality of determinations
into one identity.

Neither inference is justified without further analysis.

The plurality of determinations may concern the
internal constitution, manifestation, or relations
of an existent.
Whether such plurality also involves a plurality
of identities is a separate question.

The distinction between unity and plurality must
therefore be maintained at the level of the subject
and the respect under consideration.

% -----------------------------------------------------------------------------
\subsection{Unity Without Reduction to Homogeneity}

Unity of identity must not be confused with homogeneity.

Homogeneity concerns similarity or sameness among
elements, determinations, or features in a specified
respect.

Unity, by contrast, cannot be defined merely as the
absence of difference.

An existent may involve differentiated determinations
without those differences necessarily destroying
the unity of its identity.

Likewise, the presence of similar or homogeneous
features does not establish that several entities
possess one identity.

Similarity is not identity, and difference is not
by itself proof of distinct identity.

This distinction is important because a theory that
defines unity through homogeneity risks treating
internal differentiation as a threat to identity.

Such a theory would make the possibility of unity
depend on the reduction or elimination of difference.

The present inquiry does not adopt this assumption.

Unity and difference must be examined in relation
to one another without presuming that one necessarily
abolishes the other.

The ontological status of difference in relation
to identity is considered further in the analysis
of unity and plurality.

% -----------------------------------------------------------------------------
\subsection{The Question of the Subject of Unity}

Every claim of unity requires clarification of
the subject to which unity is attributed.

An assertion that something is unified remains
incomplete unless it specifies what is being
considered as unified.

The subject may be an existent, an identity,
a source, an aggregate, a composition,
a configuration of determinations,
or a set of manifestations.

These subjects are not interchangeable.

For example, the unity of a material aggregate
does not automatically establish the unity
of the identity of that aggregate.
Similarly, the unity of a source does not
automatically establish the unity of the identities
associated with its manifestations.

The ontological subject of unity must therefore
be distinguished from the elements, relations,
or determinations through which unity is expressed.

This requirement also applies to claims of plurality.
A plurality of components, relations, or manifestations
cannot be transferred without argument into
a claim concerning a plurality of identities.

The examination of unity must consequently specify:

\begin{enumerate}
    \item The subject whose unity is being asserted.
    \item The ontological respect in which unity is asserted.
    \item The relations or determinations relevant
    to that unity.
    \item The limits beyond which the claim of unity
    cannot be extended.
\end{enumerate}

These requirements prevent the unrestricted use
of the term unity across different ontological levels.

% -----------------------------------------------------------------------------
\subsection{The Possibility of Plurality Within an Existent}

The examination of unity of identity must leave open
the possibility that an existent may involve a plurality
of identities.

This possibility cannot be rejected merely by assuming
that the unity of an existent entails the singularity
of its identity.

At the same time, the possibility of plurality
must not be treated as an established universal fact.

The present section does not determine whether
every existent possesses one identity, whether
some existents involve several identities,
or whether the relation between an existent
and its identities admits different ontological forms.

These questions require separate investigation.

In particular, the plurality of identities must
not be inferred directly from:

\begin{itemize}
    \item The plurality of components.
    \item The plurality of manifestations.
    \item The plurality of relations.
    \item The plurality of determinations.
    \item The multiplicity of processes or transformations.
\end{itemize}

Each of these may provide a relevant condition
for further inquiry, but none constitutes sufficient
proof of plurality of identities by itself.

Conversely, the unity of an existent must not
be used as a sufficient proof that only one identity
is involved.

The question of plurality of identities is therefore
not settled by the preliminary recognition of unity.
It is examined directly in
Section~\ref{sec:plurality-of-identities}.

% -----------------------------------------------------------------------------
\subsection{The Ontological Scope of Unity of Identity}

The preceding analysis establishes several constraints
for the examination of unity of identity.

First, unity of identity is a distinct ontological
problem and cannot be reduced to the unity of
an aggregate, source, composition, or configuration.

Second, unity cannot be equated with numerical
singularity without an additional argument.

Third, the preservation of invariant determinations
does not by itself establish unity of identity.

Fourth, the plurality of determinations, components,
relations, or manifestations does not by itself
establish a plurality of identities.

Fifth, the unity of an existent does not automatically
settle the number of identities involved in its existence.

These constraints do not constitute a final theory
of unity.
They establish the distinctions required for such
a theory to be developed without confusing different
ontological questions.

The unity of identity must therefore remain open
to further analysis concerning its relation to
source, aggregate, plurality, composition,
and multiplicity.

% -----------------------------------------------------------------------------
\subsection{Conclusion: Unity as an Open Ontological Problem}

Unity of identity cannot be adequately understood
by treating unity as numerical singularity,
structural integration, homogeneity, common origin,
or the preservation of invariant features.

Each of these concerns a distinct ontological question.
Some may be relevant to the constitution or manifestation
of identity, but none can replace the analysis
of identity itself.

The unity of identity is therefore approached
as a specific ontological problem whose scope
must be established through conceptual distinctions
rather than through an assumed equivalence between
unity and other forms of oneness.

This approach also prevents the premature exclusion
of plurality of identities.
The possibility of such plurality remains open,
but it requires arguments that cannot be derived
merely from the plurality of components,
determinations, relations, or manifestations.

The following section turns to the distinction
between unity of source and unity of aggregate.
It examines two different grounds on which unity
may be attributed to an existent or to a plurality
of manifestations, without assuming that either
one determines the unity of identity.

% =============================================================================
% End of Section 02
% =============================================================================